% ============================================================================
\section{Conclusiones}\label{sec:conclusiones}
% ============================================================================

\subsection{Qué se ha hecho}

Se han implementado en C, sobre una base común de estructuras de datos, un algoritmo
constructivo, una búsqueda local con tres estructuras de entornos y cinco metaheurísticas que
siguen los esquemas de los apuntes. Sus 21 parámetros se han calibrado con un procedimiento
de un factor cada vez sobre 6 instancias, y las configuraciones resultantes se han validado
sobre otras 24. En total, 1548%
 ejecuciones, todas con solución
factible verificada de forma independiente.

\subsection{Resultados}

\begin{center}
\begin{tabular}{@{}lcl@{}}
\toprule
Algoritmo & RPD en validación & Tiempo \\ \midrule
Constructivo voraz (cociente) & 6.45\,\% & 0.01 s \\
Búsqueda local ($\D1,\D2,\D3$) & 4.67\,\% & 0.03 s \\
Algoritmo Genético & 4.28\,\% & 5 s \\
Búsqueda Tabú & 4.11\,\% & 5 s \\
Temple Simulado & 4.07\,\% & 5 s \\
VNS & 3.52\,\% & 5 s \\
GRASP & 3.45\,\% & 5 s \\
\bottomrule
\end{tabular}
\end{center}

\begin{itemize}
  \item Las cinco metaheurísticas mejoran a la búsqueda local, y la validación confirma que
        los parámetros elegidos no están sobreajustados.
  \item \textbf{GRASP (con memoria) y VNS son las mejores}, empatadas entre sí, medio punto
        por delante de temple y tabú. Las dos tienen algo en común: aplican la búsqueda
        local completa (con el entorno $\D3$) a muchas soluciones distintas.
  \item El \textbf{genético} es el más flojo; en la forma calibrada funciona, de hecho, como
        una búsqueda por perturbación (muta todos los hijos) más que como un algoritmo
        evolutivo.
  \item La ganancia de las metaheurísticas sobre la búsqueda local es modesta: entre 0.4 y
        1.2 puntos a cambio de 200 veces más tiempo.
\end{itemize}

\subsection{Por qué ninguna baja del 3.4\,\%}\label{sec:barrera}

Este es el resultado más llamativo y merece una explicación. La tabla~\ref{tab:estructura}
compara las soluciones de cada algoritmo con la solución óptima, columna a columna.

\begin{table}[H]
\centering
\caption{Estructura de las soluciones en las 6 instancias de calibración (semilla 1): RPD, número de columnas y cuántas de ellas pertenecen también a la solución óptima, que tiene 50 columnas. Medias de las 6 instancias.}
\label{tab:estructura}
\begin{tabular}{@{}lrrr@{}}
\toprule
Solución & RPD (\%) & Columnas & Comunes con la óptima \\ \midrule
Voraz, criterio 2 (mayor cobertura) & 33.43 & 49.2 & 27.2 \\
Voraz, criterio 3 (cociente) & 6.60 & 167.3 & 3.0 \\
Búsqueda local ($j_A=3$) & 4.82 & 141.7 & 3.2 \\
\midrule
Temple Simulado & 4.14 & 148.3 & 4.3 \\
Búsqueda Tabú & 4.02 & 124.5 & 4.0 \\
GRASP & 3.27 & 128.3 & 5.3 \\
Algoritmo Genético & 4.33 & 145.0 & 5.3 \\
VNS & 3.68 & 128.8 & 2.7 \\
Solución óptima & 0.00 & 50.0 & 50.0 \\
\bottomrule
\end{tabular}
\end{table}


\begin{itemize}
  \item La solución óptima usa \textbf{50 columnas grandes}. Las soluciones de todos los
        heurísticos usan entre 125 y 170 \textbf{columnas pequeñas y baratas}, y solo
        comparten con la óptima entre 3 y 5 columnas. No están ``cerca del óptimo con algún
        detalle por pulir'': son soluciones de otra clase, que casualmente cuestan casi lo
        mismo.
  \item El voraz de mayor cobertura (criterio 2) sí encuentra más de la mitad de las columnas
        óptimas (27 de 50), pero completa la solución con columnas grandes equivocadas y
        carísimas, y queda un 33\,\% por encima.
  \item Hay, por tanto, \textbf{dos familias de soluciones}: muchas columnas pequeñas (a un
        3--6\,\% del óptimo) y pocas columnas grandes (donde está el óptimo, pero también
        soluciones muy malas). Para pasar de la primera a la segunda hay que cambiar más de
        cien columnas, y los pasos intermedios son peores: tener ya la mitad de las columnas
        óptimas no garantiza nada, como muestra el criterio 2.
  \item El criterio del cociente coste/cobertura, que es el que guía la construcción, las
        reparaciones y las perturbaciones de todos los algoritmos, empuja siempre hacia la
        primera familia, porque las columnas diminutas y casi gratuitas tienen mejor
        cociente que las columnas de la solución óptima.
\end{itemize}

\begin{idea}[Interpretación]
Las cinco metaheurísticas exploran bien la familia de soluciones de columnas pequeñas (por
eso mejoran a la búsqueda local y son robustas en validación), pero ninguno de sus
mecanismos de escape (aceptar empeoramientos, prohibir movimientos, reiniciar, cruzar,
saltar) es capaz de dar el salto a la otra familia en 5 segundos. No es un problema de
calibración: la calibración mueve los resultados unas décimas, y la distancia al óptimo es
de más de tres puntos.
\end{idea}

Como se indicó en la sección~\ref{sec:optimos}, estas instancias tienen además una propiedad
especial (la relajación lineal ya da el óptimo entero) que las hace fáciles para un método
exacto. Un heurístico que usara la información de la relajación lineal o de sus variables
duales para elegir columnas tendría aquí una gran ventaja, pero eso queda fuera de las
metaheurísticas que pide la práctica.

\subsection{Limitaciones y posibles mejoras}

\begin{itemize}
  \item \textbf{Calibración de una vuelta.} No capta interacciones entre parámetros, y en
        varios parámetros ganó un valor extremo de su lista (por ejemplo \cod{alfa} del
        temple, \cod{lrc} y \cod{elite} de GRASP, o \cod{pob}, \cod{nmut} y \cod{bl} del
        genético), de modo que el mejor valor podría estar más allá. La segunda vuelta
        opcional de \cod{calibrar.py} está pensada para eso.
  \item \textbf{Tres semillas.} Suficientes para elegir parámetros con diferencias claras,
        justas para los empates. Para el informe convendría ejecutar las configuraciones
        finales con más semillas (por ejemplo 10) y aplicar pruebas estadísticas a todas las
        comparaciones.
  \item \textbf{Un único límite de tiempo.} Los parámetros son los mejores para 5 segundos.
        La figura~\ref{fig:convergencia} indica que GRASP, VNS y temple todavía mejoran al
        final, así que con más tiempo las diferencias podrían cambiar.
  \item \textbf{Una única familia de instancias.} Todas tienen el mismo tamaño y estructura;
        las conclusiones no se pueden extender sin más a otros problemas de cubrimiento.
  \item \textbf{Movimientos.} La explicación de la sección~\ref{sec:barrera} sugiere dónde
        buscar mejoras: construcciones o reparaciones que no usen solo el cociente (por
        ejemplo, alternándolo con el criterio de mayor cobertura) para sembrar soluciones de
        la otra familia.
\end{itemize}
